%% kapa: exam=1 round=31 year=2020 period=1 form=A subject=civil_law subject_ko=민법 seq=22 number=22
%% source: Q-Net 감정평가사 기출문제 (cst00309) artlSeq=5207219 file=extracted/1008/5207219/A형/A형/2020년도 제31회 감정평가사 1차 시험_1교시_A형.pdf
%% conversion: kordoc 4.24.51 -> tools/gen1.py (auto); review=proofread-sonnet; flags=none
\begin{question}{22}
부동산 물권변동에 관한 설명으로 옳지 않은 것은? (다툼이 있으면 판례에 따름)
\begin{choices}
  \item 소유권이전등기를 마친 등기명의인은 제3자에 대하여 적법한 등기원인으로 소유권을 취득한 것으로 추정되지만 그 전(前)소유자에 대하여는 그렇지 않다.
  \item 미등기건물의 원시취득자는 그 승계인과 합의하여 승계인 명의로 소유권보존등기를 하여 건물소유권을 이전할 수 있다.
  \item 등기는 물권의 존속요건이 아니므로 등기가 원인 없이 말소되더라도 그 권리는 소멸하지 않는다.
  \item 미등기건물의 소유자가 건물을 그 대지와 함께 팔고 대지에 관한 소유권이전등기를 마친 때에는 매도인에게 관습법상 법정지상권이 인정되지 않는다.
  \item 저당권설정등기가 원인 없이 말소된 때에도 그 부동산이 경매되어 매수인이 매각대금을 납부하면 원인 없이 말소된 저당권은 소멸한다.
\end{choices}
\end{question}
\answer{1}
